\section{Method}
\label{sec:method}

\paragraph{Data.}
The two-moons problem consists of two interleaving half circles with Gaussian noise added to
every point. We draw \nbval{n-train} points for training and \nbval{n-test} for testing.

\paragraph{Model.}
The classifier is a network with one hidden layer of $h = \nbval{hidden}$ $\tanh$ units and
a sigmoid output, $\nbmath{n-params}$ parameters in all, trained with the cross-entropy loss
\begin{equation}
  \mathcal{L}(\theta) = -\frac{1}{N}\sum_{i=1}^{N}
    \Big[ y_i \log p_\theta(x_i) + (1 - y_i) \log\big(1 - p_\theta(x_i)\big) \Big].
  \label{eq:loss}
\end{equation}

\paragraph{Optimisers.}
For each of the \nbval{n-seeds} seeds, all three start from the same weights and see the
same mini-batches of \nbval{batch} points for \nbval{epochs} epochs; we report the mean and
standard deviation over seeds. Writing $g_t = \nabla_\theta \mathcal{L}(\theta_t)$ for the
mini-batch gradient and $\eta$ for the learning rate, SGD takes the step
\begin{equation}
  \theta_{t+1} = \theta_t - \eta\, g_t ,
  \label{eq:sgd}
\end{equation}
momentum accumulates a velocity,
\begin{equation}
  v_{t+1} = \mu\, v_t + g_t , \qquad \theta_{t+1} = \theta_t - \eta\, v_{t+1} ,
  \label{eq:momentum}
\end{equation}
with $\mu = 0.9$, and Adam rescales each coordinate by running estimates of the first and
second moments of the gradient,
\begin{equation}
  \theta_{t+1} = \theta_t - \eta\, \frac{\hat{m}_t}{\sqrt{\hat{v}_t} + \epsilon} ,
  \label{eq:adam}
\end{equation}
with the default decay rates $\beta_1 = 0.9$ and $\beta_2 = 0.999$ of~\cite{kingma2015}.
% NOTE: the learning rates were chosen by hand, one per optimiser; they are listed in the table.
